% Working manuscript draft. Compile from this directory or the parent docs root.
% Reuses ../literature/references.bib without changing the annotated review.
% Supplemental stripe and ladder-material references are in additional_references.bib.
\documentclass[11pt,letterpaper]{article}
\usepackage[T1]{fontenc}
\usepackage{lmodern}
\usepackage[margin=1in]{geometry}
\usepackage{amsmath,amssymb}
\usepackage{microtype}
\usepackage{booktabs,tabularx,array}
\usepackage{graphicx}
\usepackage{pdflscape}
% Overleaf builds from the project root; local builds may use manuscript/.
% Keep literature/, manuscript/, and reports/ as sibling folders when uploading.
\IfFileExists{literature/references.bib}{%
  \newcommand{\projectdocroot}{.}%
  \newcommand{\manuscriptroot}{manuscript}%
}{%
  \newcommand{\projectdocroot}{..}%
  \newcommand{\manuscriptroot}{.}%
}
\usepackage{needspace}
\usepackage[numbers,sort&compress]{natbib}
\usepackage{xurl}
\usepackage{xcolor}
\usepackage[colorlinks=true,linkcolor=blue!50!black,
  citecolor=blue!50!black,urlcolor=blue!50!black]{hyperref}
\setlength{\parindent}{1.3em}
\setlength{\parskip}{0.25em}
\setlength{\emergencystretch}{3em}
\renewcommand{\arraystretch}{1.18}
\setcounter{tocdepth}{1}
\newcommand{\tp}{t_{\perp}}
\newcommand{\Ep}{E_{\mathrm p}}
\newcommand{\hc}{\mathrm{H.c.}}
\newcommand{\avg}[1]{\langle #1\rangle}
\newcommand{\ev}[1]{\hyperref[ev:#1]{Evidence~\ref*{ev:#1}}}
\newcolumntype{Y}{>{\raggedright\arraybackslash}X}

\title{Pairing, Stripes, and Competing Order\\in Coupled Hubbard Ladders\\
  \large Introduction, Background, and Current Results}
\author{Ladder MPS+MF project\\Working manuscript draft}
\date{Updated 22 September 2026}

\begin{document}
\maketitle

\begin{abstract}
Weakly coupled two-leg Hubbard ladders provide a setting in which the formation
of singlet pairs can be studied separately from the collective order selected
by coupling correlated subsystems. This draft develops the motivation for a
multichannel matrix-product-state plus mean-field study of superconducting,
charge, and magnetic order, with particular attention to the relationship
between stripe-centered pairing proposals and the broader picture of
intertwined orders. It also collects the project's current local evidence.
The isolated repulsive ladder exhibits a finite-system spin gap and substantial
pairing correlations, whereas the latest square-array trajectories at zero
nearest-neighbor interaction develop strong stripe textures. At an attractive
interaction, two contrasting initial states instead approach a common paired
state. The completed unfrustrated-cubic grid instead reaches stripes at all
nine tested coordinates. Finer square cuts retain distinct textures near
the boundary, with a small change in energy slope along the interaction cut
but no resolved first-order singularity. All four repulsive-$V$ square
starts lose pairing, including both intertwined seeds. Both trellis seeds
become paired in a skew one-ladder cell but develop stripes in a rectangular
two-ladder cell, where large alternating relaxation persists. Complete
retrospective measurements show that these striped endpoints retain
short-distance singlet-pair correlations concentrated near hole-rich magnetic
walls, while longer-distance correlations are strongly suppressed.
Rung--leg signs predominantly survive, but reference and fit-window
sensitivity preclude a single asymptotic decay exponent. All four newer
square two-ladder runs at the attractive paired coordinates pass fixed-point
gates and reproduce the paired state from both tested seeds. These
observations sharpen the distinction between local pairing
tendencies and the order selected in an array. The numerical results remain
finite-length, finite-bond-dimension evidence, with explicit limits on
self-consistency, response completeness, and thermodynamic interpretation.
\end{abstract}

\noindent\textbf{Draft scope.}
Sections~\ref{sec:intro}--\ref{sec:background} are draft manuscript prose.
Section~\ref{sec:results} is a dated collection of results and interpretation,
including unresolved points that must remain visible during manuscript
development. The evidence cutoff is the synchronized analyses through
22 September 2026; the literature discussion incorporates the stripe and
ladder-material discussions through 18 September. Appendix~\ref{sec:evidence} identifies the local
sources. The new analysis reads completed measurements and simulations;
no additional many-body solve was performed for manuscript preparation.

\tableofcontents
\clearpage

\section{Introduction}
\label{sec:intro}

The relationship between superconductivity, antiferromagnetism, and spatial
charge organization is a central question in doped Mott insulators.
The Hubbard model provides a minimal language for studying how kinetic energy
and local repulsion generate collective behavior, while its extensions allow
the influence of nonlocal interactions and lattice structure to be examined
systematically~\citep{Hubbard1963,Anderson1987,Lee2006,Arovas2022}.
In this setting, strong singlet correlations need not lead directly to
superconductivity. Charge and spin order can develop from the same interacting
electrons, and the scales associated with pair formation and collective phase
coherence can differ. A microscopic account of superconductivity must therefore
address the competing ways in which a correlated electronic state can organize.

Stripe order makes this issue particularly concrete. Experiments in
lanthanum-based cuprates established charge modulations accompanied by
incommensurate magnetic order, with antiferromagnetic domains separated by
hole-rich regions~\citep{Tranquada1995}. Subsequent measurements found closely
related magnetic excitation spectra in stripe-ordered and superconducting
materials, motivating theories in which stripe correlations participate in
the pairing mechanism~\citep{Tranquada2004}. A complementary view regards
stripes as a competing organization of a doped Mott insulator whose
short-distance singlet physics can also support superconductivity without
requiring a particular stripe arrangement~\citep{Lee2006}. The resulting
question concerns both the origin of pairing and the consequences of charge
organization for the coherence of pairs.

The distinction between these views is sharpened by separating the existence
of stripe correlations from their necessity for superconductivity. In
stripe-ordered La$_{1.875}$Ba$_{0.125}$CuO$_4$, transport and magnetization
measurements reveal superconducting correlations within the planes despite
strong suppression of three-dimensional phase
coherence~\citep{Li2007}. A spatially sign-changing superconducting order
parameter can frustrate interplane Josephson coupling and offers a specific
explanation of this separation~\citep{Berg2007}. Such observations permit
cooperation at the level of local pairing and competition at the level of
collective order. They motivate treating superconducting, charge, and spin
degrees of freedom together, as emphasized in the framework of intertwined
orders~\citep{Fradkin2015,Agterberg2020}.

Modern Hubbard-model calculations reinforce the dependence of this
relationship on the microscopic Hamiltonian. Complementary numerical methods
find stripe order with small energy differences between wavelengths in a
studied regime of the nearest-neighbor square-lattice
model~\citep{Zheng2017}. Other calculations find no superconducting ground state
in specified regimes of that minimal model~\citep{Qin2020}, while additional
hopping terms can support pairing together with spatial charge
modulation~\citep{JiangDevereaux2019,Xu2024}. Deliberately imposed hopping
inhomogeneity can also enhance superconducting correlations on
cylinders~\citep{JiangKivelson2022}. These results do not give a universal sign
to the influence of stripes. They instead identify a competition that must be
resolved at the level of a specified model, geometry, and numerical accuracy.

Two-leg ladders provide useful correlated building blocks for this problem.
Their magnetic spectrum differs qualitatively from that of a single chain,
and doped repulsive ladders can support a spin gap, bound holes, and slowly
decaying singlet-pair correlations~\citep{Dagotto1992,Noack1994,DagottoRice1996}.
The corresponding low-energy description can contain a gapped spin sector and
a gapless total-charge mode~\citep{LutherEmery1974,BalentsFisher1996}.
High-accuracy numerical studies show that the relative decay of pair and
density correlations requires careful control of open boundaries, bond
dimension, and fit windows~\citep{Dolfi2015,Shen2023}. Thus a ladder can exhibit
substantial pairing tendencies while retaining the fluctuations characteristic
of an effectively one-dimensional system.

Cuprate ladder materials make the role of transverse geometry an experimental
question. Superconductivity was discovered under pressure in
Sr$_{0.4}$Ca$_{13.6}$Cu$_{24}$O$_{41.84}$, a member of the
Sr$_{14-x}$Ca$_x$Cu$_{24}$O$_{41}$ chain--ladder family~\citep{Uehara1996}.
The same family exhibits ladder charge order and pressure-dependent transport,
providing a setting in which pairing, charge organization, and interladder
coherence can be compared~\citep{Nagata1998,Abbamonte2004,Rusydi2006}.
Its ladder sheets have displaced neighbors connected by zigzag paths: a
trellis network rather than an array with only aligned transverse bonds.
The chain subsystem also acts as a charge reservoir, so varying composition
or pressure changes more than the coupling between otherwise fixed ladders.
These distinctions motivate a controlled comparison of transverse
connectivities before attempting a material-specific phase diagram.

Coupling such ladders introduces a second question: which order develops when
the subsystems reinforce one another? Stripe-centered theories provide an
explicit route from spin-gap formation to pairing and then to Josephson
coherence~\citep{Emery1997}. A spatially modulated repulsive Hubbard
construction demonstrates how nearly decoupled ladders can supply a pairing
scale while weak interladder coupling establishes superconducting
order~\citep{Arrigoni2004}. Matrix-product-state plus mean-field methods
provide a practical numerical realization of the broader coupled-subsystem
strategy~\citep{Bollmark2020,Bollmark2023}. Their multichannel extension shows
why charge and superconducting solutions must be retained and compared within
a common framework~\citep{Bollmark2025}.

The interaction between neighboring densities supplies an additional tuning
parameter. Calculations on extended Hubbard ladders and cylinders find that
attraction can enhance pairing over selected parameter ranges, while strong
attraction can also favor charge segregation~\citep{Zhou2023,Peng2023}.
Recent cuprate-ladder spectroscopy, interpreted together with DMRG, provides
experimental motivation for considering attraction beyond the simple onsite
Hubbard interaction~\citep{Padma2025}. Its implications for a coupled array
remain to be determined: stronger binding changes both the properties of the
individual ladder and the effective interladder couplings. Pair mobility,
susceptibility, and competing order must be considered along with binding.

The present project investigates these issues in arrays of repulsive two-leg
extended Hubbard ladders, treating the ladder with an MPS ground-state method
and the transverse environment with a multichannel mean-field construction.
The principal regime is $U/t=8$ and density $n=15/16$ per physical site.
Rung hopping, nearest-neighbor interaction, and transverse connectivity are
varied while pairing, normal exchange, and spin-resolved density fields remain
available. In particular, the trellis comparison tests whether zigzag
coupling changes the balance between static stripes and pairing at the same
intraladder parameters, using both one- and two-ladder spatial cells.
The immediate objective is to establish reproducible ordered
solutions and understand how their competition relates to the underlying
ladder correlations. The current evidence suggests that intrinsic pairing
tendencies can coexist with an array tendency toward static stripes, and that
attraction can reverse the observed outcome at a representative square-array
point. Establishing the extent and thermodynamic stability of these behaviors
requires the numerical controls discussed below.

\section{Physical background and connection to the model}
\label{sec:background}

\subsection{Local singlets, pairing correlations, and coherence}

Three observables answer different questions. A spin gap measures the energy
cost of the lowest accessible spin excitation at fixed particle number.
A pair-binding energy compares even- and odd-particle sectors. A pair
correlation function measures how a pair created in one region correlates with
a pair removed elsewhere. None of these quantities alone establishes a
nonzero superfluid stiffness or a superconducting transition temperature.
For example, an isolated number-conserving ladder has zero anomalous
expectation value even when its pair correlations decay algebraically.

For a spin-singlet bond operator,
\begin{equation}
 \Delta_{ab}=\frac{c_{a\uparrow}c_{b\downarrow}
                   -c_{a\downarrow}c_{b\uparrow}}{\sqrt{2}},
 \qquad
 P_{ab,cd}=\avg{\Delta^\dagger_{ab}\Delta_{cd}},
 \label{eq:pair-op}
\end{equation}
the relative signs between different bond orientations characterize the
internal pair structure. Opposite signs on leg and rung bonds are commonly
described as $d$-wave-like in a two-leg ladder. The reduced lattice symmetry
and discrete transverse momentum should be kept in mind when connecting this
terminology to the gap of a two-dimensional cuprate.

An order parameter with a nonzero center-of-mass wavevector is a different
property. Schematically,
\begin{equation}
 \avg{\Delta_b(R)}=f_b\left[\Psi_0+
          \Psi_Q\cos(QR+\phi)\right],
 \label{eq:pdw}
\end{equation}
separates the internal bond form factor $f_b$ from the spatial structure of
the condensate. An internal leg/rung sign change alone is not evidence of a
pair-density wave. Nor does a boundary-induced modulation of a predominantly
uniform pair field establish a distinct finite-momentum superconducting
phase~\citep{Agterberg2020}.

\subsection{Two views of stripes and their modern interpretation}

In a common-Mott-physics interpretation, singlet and magnetic correlations
provide a shared microscopic starting point for several possible ordered
states. Static stripes are then one possible competitor, and their presence
need not be a prerequisite for pairing~\citep{Lee2006}. In a stripe-centered
interpretation, spatial organization helps create the conditions for pairing.
The spin-gap proximity mechanism, for example, couples a conducting stripe
to a magnetically correlated environment; pair hopping transmits a spin gap
to the mobile carriers, while Josephson coupling acts at a separate stage to
establish coherence~\citep{Emery1997}.

These interpretations share an emphasis on strong correlations and singlets.
Their sharper disagreement concerns whether stripe organization is an
essential intermediate step. The observations of planar superconductivity
in stripe-ordered LBCO and the associated layer-decoupling proposal show how
static stripes can coexist with pairing while suppressing a particular route
to bulk coherence~\citep{Li2007,Berg2007}. Conversely, coexistence in a model
or material does not establish that stripes caused its pairing.

The relation to the pseudogap remains unresolved. NMR measurements in Eu-LSCO
find spin-stripe order extending practically to the pseudogap endpoint and
identify competition with superconductivity, providing evidence for a close
connection without determining its causal origin~\citep{Missiaen2025}.
The spin gap of an isolated ladder, an electronic pseudogap in a cuprate, and
the superconducting quasiparticle gap are different observables. A calculation
of one should not be presented as a measurement or explanation of the others.

\subsection{What magnetic spectroscopy establishes}

Incommensurate magnetic peaks can be understood through antiphase spin
domains. With lattice spacing set to one, a schematic stripe texture is
\begin{equation}
 \avg{S^z_{x,y}}=m(-1)^{x+y}\cos(\kappa x),
 \qquad q_{\mathrm c}=2\kappa,\qquad
 \mathbf q_{\mathrm s}=(\pi\pm\kappa,\pi).
 \label{eq:stripe-cartoon}
\end{equation}
The charge period is half the period of the staggered spin envelope. This
spatial relation is useful evidence for an antiphase stripe texture, but its
presence on a finite open ladder does not determine the thermodynamic
amplitude or preferred wavelength.

The similarity between magnetic spectra in stripe-ordered LBCO and YBCO was
an important motivation for stripe-centered
interpretations~\citep{Tranquada2004}. Quantum stripe calculations show that
low-energy spin-wave-like excitations and higher-energy ladder-like triplet
excitations can occur within the same ordered model~\citep{Vojta2004}.
Consequently, a failure of a simple semiclassical intensity pattern does not
rule out stripes or all low-energy spin waves. The connection between a
spectral shape and a microscopic pairing mechanism requires more than
qualitative resemblance.

The magnetic ladders in such spin-only constructions represent magnetic
regions of a stripe pattern. The present model instead uses doped electronic
ladders containing mobile charge. Moreover, its current analyses concern
ground-state profiles and equal-time correlations, rather than a computed
$S(\mathbf q,\omega)$. They can test the static antiphase relation in
Eq.~\eqref{eq:stripe-cartoon}, but cannot yet demonstrate an hourglass
spectrum or identify a neutron resonance.

\subsection{Extended Hubbard ladders, materials, and transverse coupling}
\label{sec:ladder-materials}

Magnetic ladder insulators establish the relevance of the ladder building
block, but do not by themselves realize a doped electronic array.
The spin gap observed in SrCu$_2$O$_3$ is an early example~\citep{Azuma1994}.
Introducing mobile holes is a separate materials problem: chemical
substitution can modify exchange paths, disorder, and carrier localization.
Hole doping is possible in La$_{1-x}$Sr$_x$CuO$_{2.5}$, where an
insulator-to-metal transition was observed without superconductivity down
to 5~K in the reported measurements~\citep{Hiroi1995}.
Its coupled-ladder electronic structure has substantial three-dimensional
coupling~\citep{NormandRice1996}. Such compounds motivate studying coupled
ladders and dimensionality, but do not identify our particular square or
cubic kernels with a complete material Hamiltonian.

The superconducting chain--ladder cuprates provide a more direct motivation
for the trellis extension. Uehara et al.\ reported superconducting onset
temperatures of 12~K at 3~GPa and 9~K at 4.5~GPa in
Sr$_{0.4}$Ca$_{13.6}$Cu$_{24}$O$_{41.84}$~\citep{Uehara1996}.
The structure contains both Cu$_2$O$_3$ two-leg ladder sheets and CuO$_2$
chains. Neighboring ladders are offset longitudinally, producing the two
zigzag paths of a trellis network~\citep{Popovic2000}.
Ca substitution is isovalent; it redistributes holes between chains and
ladders rather than simply adding a known number of ladder holes.
NQR and pressure-dependent optical measurements likewise support
pressure-induced transfer of holes into the ladders~\citep{Fujiwara2009,Frank2014}.
Consequently, the model density $n=15/16$ cannot be assigned directly to a
Ca concentration or experimental pressure.

The experimental charge-order landscape is richer than a monotonic conversion
from a CDW to a superconductor. Transport, dielectric, and optical measurements
found a CDW transition scale decreasing from about 210~K at $x=0$ to 10~K at
$x=9$~\citep{Vuletic2003}. Resonant x-ray scattering established a ladder
hole crystal in Sr$_{14}$Cu$_{24}$O$_{41}$ and commensurate charge order also
at $x=10,11,12$~\citep{Abbamonte2004,Rusydi2006}. High Ca content alone
therefore does not establish the absence of CDW order. Composition, pressure,
and commensurability must be specified when comparing a uniform paired model
state with experiment.

Static magnetic order and dynamic spin correlations must also be separated.
Neutron measurements on Sr$_{2.5}$Ca$_{11.5}$Cu$_{24}$O$_{41}$ found that the
ordered moment decreased under pressure and appeared to disappear near the
superconducting threshold~\citep{Katano1999}. A paired state without a
resolved SDW is compatible with that trend, but vanishing static magnetization
does not determine the spin gap. High-pressure NMR found gapped quasiparticle
excitations in the superconducting state, described as $s$-wave-like in that
specific sense~\citep{Fujiwara2003}. This is not a phase-sensitive measurement
of the relative signs of leg and rung pairing. A ladder's $d$-wave-like bond
structure need not force nodes on its Fermi surfaces; neither agreement nor
disagreement of pairing symmetry follows from our bond signs alone.
RIXS on Ca-substituted material further motivates including localization:
the observed magnetic spectra were interpreted using off-ladder impurity
potentials that immobilize holes~\citep{Tseng2022}.

Each ladder has $L$ rungs, two legs $\ell=0,1$, and $2L$ physical sites.
Its Hamiltonian is
\begin{align}
 H_{\mathrm{lad}}={}&
 -t\sum_{x=1}^{L-1}\sum_{\ell,\sigma}
  (c^\dagger_{x\ell\sigma}c_{x+1,\ell\sigma}+\hc)
 -t_0\sum_{x,\sigma}
  (c^\dagger_{x0\sigma}c_{x1\sigma}+\hc)\nonumber\\
 &+U\sum_{x,\ell}n_{x\ell\uparrow}n_{x\ell\downarrow}
 +V\sum_{\avg{a,b}_{\mathrm{lad}}}n_a n_b .
 \label{eq:model}
\end{align}
The final sum includes nearest-neighbor leg and rung bonds.
Thus $V<0$ is attractive in the density channel. Open boundaries are used
along the ladder; $n=N/(2L)$ is the density per physical site and
$p_{\mathrm h}=1-n$ the hole density. The rung hopping $t_0$ and interladder
hopping $\tp$ must not be confused with the next-nearest-neighbor hopping
$t'$ in two-dimensional Hubbard studies.

An approximate parameter scale comes from the optical analysis of
Sr$_{14}$Cu$_{24}$O$_{41}$ by Popovi\'c et al.~\citep{Popovic2000}:
$t=t_0=0.26$~eV, interladder hopping $t_{xy}=0.026$~eV, and $U=2.1$~eV.
Interpreting $t_{xy}$ as the amplitude on each equivalent zigzag path gives
$U/t\simeq8.1$, $t_0/t=1$, and $\tau_0/t=\tau_1/t=0.1$, motivating our
$(8,1,0.1,0.1)$ comparison. These are effective parameters from an optical
model of the parent compound, not a pressure-resolved fit to Uehara's
superconducting sample. That analysis does not independently determine the
nearest-neighbor $V$ used here; $V=0$ is our baseline assumption.

There is independent motivation to vary $V$. The RIXS--DMRG comparison of
Padma et al.\ favors an effective attraction $V/t\simeq-1$ to $-1.25$ for
Sr$_{14}$Cu$_{24}$O$_{41}$~\citep{Padma2025}. Their model uses
$t=0.38$~eV, $t_0/t=0.84$, and an additional \emph{intraladder} diagonal
hopping $t'/t=-0.3$, distinct from our interladder $\tau_0,\tau_1$.
It adjusts $U$ with $V$ to hold the exchange scale fixed, whereas our scans
hold $U/t=8$. Thus its fitted attraction motivates extended interactions
without providing a directly transferable parameter set for the present
Hamiltonian. Pairing at $V=0$ in our approximation would not establish that
the material has no effective attraction. Chains, explicit oxygen orbitals,
structural modulation, disorder, and longer-range Coulomb interactions are
also absent from Eq.~\eqref{eq:model}.

The transverse reduction uses the positive scale
\begin{equation}
 g=\frac{\tp^2}{|\Ep|},\qquad
 \Ep^{(h)}=E(N-2,S_z=0)+E(N,S_z=0)
                -2E(N-1,S_z=\tfrac12),
 \label{eq:kernel}
\end{equation}
where negative $\Ep^{(h)}$ denotes favorable hole binding. Geometry-dependent
kernels convert anomalous, exchange, and density correlators into fields for
the next ladder solve. The density is targeted through a chemical potential.
The present ansatz retains real, spin-diagonal normal fields and real
anomalous fields; it does not provide unrestricted complex phases or
noncollinear magnetism (\ev{methods}).

For illustration, define $\gamma=2g$. The same-leg pairing kernels, organized
by the two leg indices, are
\begin{equation}
 K_{\alpha}^{\mathrm{square}}=\gamma
 \begin{pmatrix}0&1\\1&0\end{pmatrix},\quad
 K_{\alpha}^{\mathrm{cubic,u}}=\gamma
 \begin{pmatrix}0&3\\3&0\end{pmatrix},\quad
 K_{\alpha}^{\mathrm{cubic,f}}=\gamma
 \begin{pmatrix}2&1\\1&2\end{pmatrix}.
 \label{eq:geometries}
\end{equation}
These matrices summarize one part of the full map, not the entire
Hamiltonian. An equation-level comparison identifies the repulsive-ladder
benchmark of Ref.~\citep{Bollmark2023} with the branch called
\emph{cubic frustrated} in this project. That name is the project's
nomenclature. Its published result should not be assigned to all three
geometries (\ev{geometry}).

The trellis adds a longitudinal displacement to the transverse map.
Writing $m$ for the ladder index, a minimal hopping Hamiltonian is
\begin{equation}
 H_{\mathrm{zigzag}}=-\sum_{m,i,\sigma}
 c^\dagger_{m,i,1,\sigma}
 \left(\tau_0 c_{m+1,i,0,\sigma}
       +\tau_1 c_{m+1,i+1,0,\sigma}\right)+\hc .
 \label{eq:trellis-hopping}
\end{equation}
The symmetric case has $\tau_0=\tau_1$; removing the shifted path recovers
square connectivity in the one-ladder embedding. Triangular connections can
frustrate interladder antiferromagnetic alignment, making this geometry a
natural test of whether static stripe order remains favored over pairing.
The outcome must be calculated with both channels present; it is not implied
by frustration alone.

In the same common-denominator approximation, let $T_{mn}$ map sites on
neighboring ladder $n$ to ladder $m$, including both paths, and define
$F_{ab}=\avg{c_{b\uparrow}c_{a\downarrow}}$ and
$X^\sigma_{ab}=\avg{c^\dagger_{a\sigma}c_{b\sigma}}$.
For the real fields retained here, the kernels take the form
\begin{align}
 \alpha_m&=\frac{2}{|\Ep|}\sum_{n\sim m}
 \mathcal P\!\left[T_{mn}(\mathcal P F_n)T_{mn}^{T}\right],\\
 B_{m\sigma}&=\frac{2}{|\Ep|}\sum_{n\sim m}
 \mathcal P\!\left[T_{mn}\mathcal P
 (X_n^\sigma-\tfrac12 I)T_{mn}^{T}\right].
 \label{eq:trellis-fields}
\end{align}
Here $\mathcal P$ retains rung separations within the chosen range cutoff,
and $B$ combines diagonal density and off-diagonal normal fields. Each
neighbor contributes separately; the two hops of a pair transfer are not
split between different neighboring ladders. The shifted paths introduce
mixed $\tau_0\tau_1$ terms, so a coordination-number rescaling of
Eq.~\eqref{eq:geometries} is insufficient. In particular, diagonal normal
fields can depend on bond correlations as well as density (\ev{trellis}).

Reciprocity fixes the apparent directional asymmetry of the rung indexing.
With $A=\tau_0 I+\tau_1 S$ and $S_{ij}=\delta_{j,i+1}$, the one-ladder
map contains $A F^{00}A^T$ on leg 1 and $A^T F^{11}A$ on leg 0, with the
range projections understood. Opposite shifts belong to opposite spatial
neighbors at every iteration; alternating the indexing convention between
iterations would instead alternate the map being solved. The comparison
therefore uses a fixed reciprocal one-ladder cell and an explicit rectangular
two-ladder $A/B$ cell. Both are stationary spatial ansatzes, with different
allowed transverse arrangements of longitudinal textures and different open
end cuts. The two-ladder calculation targets the same mean density on each
ladder separately; it does not allow unequal average fillings (\ev{trellis}).

These runs test whether geometry can change the selected order at fixed
intraladder parameters and $V=0$. A paired solution without resolved bulk
CDW/SDW would be qualitatively compatible with a superconducting regime of
the chain--ladder family, but would not locate it in pressure or reproduce
the material's whole phase diagram. The first available trellis trajectory
and its unresolved convergence and cell-comparison checks are described in
Section~\ref{sec:positive-trellis}. Its anomalous amplitude is not a
measurement of $T_c$, phase stiffness, or a dynamical spin gap.

This construction is conceptually related to coupled-ladder pairing
mechanisms but does not implement every assumption of the original stripe
proposals. In particular, the explicit alternating dopings in the period-four
construction of Ref.~\citep{Arrigoni2004} differ from an array of equivalent
ladders, although that work also considers an equivalent-ladder period-two
limit.

We study equivalent coupled ladders as a prescribed microscopic geometry.
The observed stripe texture is an additional longitudinal charge and spin
modulation; within the repeated-ladder embedding, its domain walls extend
transverse to the ladder axes. This transverse alignment is a restriction
of the embedding rather than an independently optimized arrangement of
distinct ladders. Collective transverse quantum fluctuations are
approximated by static mean fields.

\section{Results and interpretation accumulated to date}
\label{sec:results}

\subsection{Evidence hierarchy and common parameters}

The strongest current comparison concerns the square geometry at
$t_0/t=1.4$, $\tp/t=0.1$, $U/t=8$, and $n=15/16$, with $V/t=0$ and
$-0.4$. The latest paired and striped starts use $L=64$ rungs and bond
dimension $\chi=200$. A separate $V=0$ comparison tests $\chi=400$, and
the isolated-ladder backbone reaches $\chi=1200$. These are different
calculations with different error controls, rather than one uniformly
converged data set.

It is useful to distinguish an observed texture, a numerically
self-consistent state, stability against specified perturbations, and a
thermodynamic ground state. The recent trajectories provide strong texture
and basin information. Their remaining acceptance failures and the incomplete
length and bond-dimension studies prevent the last two categories from being
identified without qualification. Here, \emph{basin} describes access under a
specified numerical map, not a measured free-energy basin.

\subsection{The isolated ladder already has substantial pairing tendencies}
\label{sec:bare}

The six-sector backbone at $V=0$ and $t_0/t=1.4$ gives the finite-system
quantities in Table~\ref{tab:bare}. All six final $\chi=1200$ sectors pass
their stored numerical gates. Changes from $\chi=800$ are useful sensitivity
observations, but not all sectors pass at that earlier stage; the comparison
therefore does not constitute a complete bond-dimension extrapolation
(\ev{bare}, \ev{review}).

\begin{table}[htbp]
\centering
\caption{Isolated-ladder evidence at $L=64$, $U/t=8$, $n=15/16$,
$t_0/t=1.4$, and $V=0$, from the final $\chi=1200$ backbone.
Negative binding energies use Eq.~\eqref{eq:kernel}'s convention.}
\label{tab:bare}
\begin{tabularx}{\textwidth}{l r Y}
\toprule
Quantity & Value & Interpretation \\
\midrule
Spin gap $\Delta_s/t$ & $0.1533613$ & Fixed-particle-number spin excitation \\
Even-particle charge gap $\Delta_{c,2}/t$ & $0.01085830$ & Finite-size charge curvature \\
Hole binding $\Ep^{(h)}/t$ & $-0.1465102$ & Favorable two-hole binding \\
Particle binding $\Ep^{(e)}/t$ & $-0.1485427$ & Favorable two-particle binding \\
Pair decay exponent & $0.7614\pm0.0215$ & Fit to rung-pair correlations \\
Charge decay exponent & $1.2158\pm0.0830$ & Less reliable fit to charge correlations \\
\bottomrule
\end{tabularx}
\end{table}

The measured pair correlations decay more slowly than the fitted charge
correlations. Pair fits have $R^2$ between $0.953$ and $0.982$, while charge
fits give $0.527$--$0.707$. The quoted fit errors do not include the full
finite-size, boundary, and truncation uncertainty. These results are
consistent with a spin-gapped ladder with appreciable pairing correlations,
but do not by themselves establish its asymptotic exponents or the ordered
state of an array~\citep{Dolfi2015,Shen2023}.

Two other diagnostics illustrate the importance of conventions and fit
quality. The stored charge parameter is $1.4247$ with rung normalization and
$0.7124$ with site normalization. An unqualified symbol $K_\rho$ would
therefore be ambiguous. A central-charge fit yields approximately $2.009$
with $R^2=0.155$ and is excluded from physical interpretation. The
open-boundary density profile has a strong component at $q/\pi=1/8$,
which can arise as a Friedel response and is not independent proof of
spontaneous charge order~\citep{WhiteAffleckScalapino2002}.

The Stage~1 equal-time covariance spectra are broad, with participation
ranks around $48$--$63$. The first bulk spin-odd candidate follows two
edge-localized modes, and a charge mode at $q/\pi=1/8$ first appears at rank
11. This motivated combining physically chosen fields with covariance-based
directions instead of assuming that a handful of leading covariance modes
would exhaust the relevant response. Equal-time covariance is itself distinct
from a static susceptibility, which also depends on excitation energies.

\subsection{Response discovery: a strong pairing direction with unresolved controls}
\label{sec:response}

The Stage~2 pilot reduces fourteen named fields to twelve independent probe
directions: nine normal and three pairing directions. Its projected
mean-field response has a leading square pairing multiplier of about
$4.6688$ and a leading positive normal multiplier of $1.1879$.
The cubic cases both give a leading pairing multiplier of $24.4730$,
while their leading normal behavior differs: the largest-magnitude normal
multiplier is approximately $-1.1879$ for the frustrated construction and
$+3.5637$ for the unfrustrated one (\ev{response}).

These values show that the isolated ladder supplies strong candidate pairing
response and that transverse geometry changes the competing normal channels.
They do not establish which nonlinear state wins. The discovery reciprocity
error is approximately $0.121\%$ and the normal/pair cross-block ratio
$0.332\%$, but the maximum omitted-response norm fraction is $0.83267$.
Small errors within the projected matrix can therefore coexist with a large
response outside the chosen subspace. The second-amplitude validation was not
present in the reviewed evidence (\ev{review}).

More fundamentally, the zero-field isolated ladder is not a fixed point of
the full normal mean-field map: its square-map image has a normalized field
norm of approximately $0.0306655$, dominated by normal fields. A derivative
at that reference is a bare-ladder diagnostic, not the stability matrix of a
self-consistent dressed normal state. The associated square estimate
$\tp/t\simeq0.04628$ must consequently remain a projected bare-reference
onset estimate, not a phase boundary.

This resolves the apparent tension between a large bare pairing response and
the later nonlinear stripe trajectories. They concern different reference
states, finite-amplitude feedback, and restricted response information.
Positive or negative multipliers describe growth or alternation of the
specified iteration. Neither is a physical excitation frequency or a global
energy criterion. A physical two-sublattice interpretation of a raw two-cycle
also requires a mapping to the transverse lattice and the appropriate
partner-field energy~\citep{Bollmark2025}.

\subsection{Why the earlier paired-state interpretation was revised}

The loose six-seed $V=0$ campaign initially produced six software-accepted,
pairing-bearing endpoints, with a corrected energy spread of only
$6.0274\times10^{-7}t$ per physical site. This suggested a common paired
region, but the later history analysis showed that it did not establish
a unique stable paired basin. Three histories stopped after six evaluations;
three others showed coherent spin-density growth during raw iteration and
subsequent suppression under Anderson acceleration. The conspicuous change
in one trajectory is reproduced by its first Anderson coefficients,
$+22.6811$ and $-21.6811$ (\ev{basins}).

This finding does not imply that Anderson acceleration changes the
fixed-point equation or that every accelerated solution is unphysical.
It shows that short histories and accelerated access do not test stability
against a resolved competing-order perturbation. The subsequent 95\%/5\%
two-family experiment addresses that question by using representative paired
and striped reference fields, adding a finite component of the competing
order, and retaining raw updates throughout.

A related warning comes from the square $(t_0/t,V/t)=(1.0,-0.4)$ coverage
point. Its \emph{diverging} outcome is a convergence classification: the fields
remain bounded while pairing gives way to a stripe-like texture and an
accelerated residual spike triggers the stop. It is not evidence of a
divergent physical observable or an accepted phase (\ev{grid}).

\subsection{Attraction: two contrasting starts approach a common paired state}
\label{sec:attraction}

At $(t_0/t,V/t)=(1.4,-0.4)$, both original 80-evaluation trajectories approach
the same pairing-dominated state. The initially strong stripe's bulk
leg-odd spin Hartree RMS decreases from about $1.99\times10^{-2}t$ to
$3.88\times10^{-8}t$; the pairing start ends near $2.49\times10^{-8}t$.
Both pairing-field RMS values settle near $2.8236\times10^{-3}t$.
Bulk RMS here uses rungs 6--59 (\ev{attractive}).

The physical symmetrized anomalous correlators have bulk leg amplitudes near
$+0.035217$ and rung amplitudes near $-0.053242$, establishing a
$d$-wave-like internal sign pattern in the stored convention. These are
not normalized versions of Eq.~\eqref{eq:pair-op} or pair--pair correlators.
The full anomalous matrices agree to approximately $5.85\times10^{-5}$
in relative L2 norm, and the largest site-density difference is below
$6.23\times10^{-6}$. The residual boundary density profile remains visible;
its survival does not indicate a surviving strong static spin stripe.

The terminal corrected energies differ by only
$7.49\times10^{-9}t$ per site, supporting agreement of the two trajectories
at their resolved accuracy. Both historical acceptance flags remain false.
The principal issues are residual fluctuations in nearly vanishing channels,
an oversensitive scalar slow-mode check, and a marginal energy-window miss
in one lineage. The later noise-aware history checks support the plateau
interpretation, but do not retroactively certify every original gate.
Thus the result is strong evidence for a common paired attractor at these
numerical settings, without proof of a global minimum or higher-$\chi$
stability.

\subsection{Zero attraction: stripe order grows while anomalous pairing recedes}
\label{sec:zero}

The corresponding $V=0$ trajectories behave differently. The stripe start
reaches 80 evaluations with bulk leg-pairing field RMS
$1.75\times10^{-10}t$ and spin Hartree RMS $0.02414t$.
The pairing start ends at the wall-time limit after 62 evaluations:
its pairing RMS has fallen from $0.004502t$ to $0.0007090t$, while spin
RMS has grown from $0.001816t$ to $0.02251t$. Pairing falls by another
$65.6\%$ across the final five records. The trend is already present before
the deadline-limited final density solve (\ev{zero}).

\begin{figure}[p]
\centering
\includegraphics[width=\textwidth]{\projectdocroot/reports/two_basin_v000_20260912/profile_evolution.pdf}
\caption{Previously generated spatial-history figure for the two $V=0$
anchor trajectories. The upper row shows the staggered leg-odd spin
mean field; the lower row shows the leg-even pairing mean field.
The two lower panels use different vertical scales. The stripe start loses
pairing, while the pairing start develops a stripe texture and its pairing
decreases. Curves are self-consistency evaluations, not physical time
evolution. The final states remain unaccepted for resolved field drift.
Source: \ev{zero}.}
\label{fig:zero}
\end{figure}

The residual anomalous correlators in the pairing-start trajectory still
have opposite leg/rung signs, but their amplitudes are decreasing rapidly.
They do not establish settled coexistence. The two spin textures are similar
but not identical: the bulk physical spin profiles differ by $26.4\%$ in
relative L2 norm, and the density difference can reach $0.04365$ per site.
The terminal global relative residuals are about $5.88\times10^{-4}$ and
$2.65\times10^{-2}$ for the stripe and pairing starts, respectively,
against a $10^{-4}$ threshold.

The working interpretation is that a robust stripe texture has been
identified and that the finite competing-order perturbation drives the
pairing-dominated start toward it under raw iteration. Strict spatial
self-consistency remains pending. In particular, the stripe amplitude's
apparent saturation does not eliminate slow rearrangement of the profile.
These observations support strong competition between static spin/charge
order and the anomalous pairing solution at this point. They do not prove
that the final stripe is the global minimum or that all local pairing
correlations have disappeared.

\subsection{Stripe wavevectors and slow texture relaxation}

Both $V=0$ endpoint profiles have dominant full-length Fourier modes
$m_{\mathrm c}=4$ for charge and $m_{\mathrm s}=30$ for leg-odd spin.
With $q_m=2\pi m/L$ and $L=64$,
\begin{equation}
 q_{\mathrm c}=\frac{\pi}{8},\qquad
 q_{\mathrm s}=\frac{15\pi}{16},\qquad
 \kappa=\pi-q_{\mathrm s}=\frac{\pi}{16}.
 \label{eq:observed-q}
\end{equation}
The observed relation $q_{\mathrm c}=2\kappa$ and the sign reversals of the
staggered spin envelope near charge-depleted regions support an antiphase
stripe interpretation. The charge period is 16 rungs and the spin-envelope
period 32 rungs.

At hole density $p_{\mathrm h}=1/16$, the magnetic displacement in
reciprocal-lattice units is $\delta_{\mathrm s}=\kappa/(2\pi)=1/32
=p_{\mathrm h}/2$. It differs from the often discussed cuprate relation
$\delta_{\mathrm s}\simeq p_{\mathrm h}$. In an ideal domain-wall counting
picture with stripe filling $\nu$ holes per unit transverse length,
$p_{\mathrm h}=\nu/d$ and $\delta_{\mathrm s}=1/(2d)$; the present
wavevectors would correspond to $\nu=1$, rather than the $\nu=1/2$ used
in the usual half-filled-stripe cartoon. This is a conditional interpretation
of the finite-system wavelength, not a measurement of wall filling.
Neither the preferred wavelength nor its scaling with doping has been
established.

The September~13 follow-up explains why the stripe looks settled while
failing strict convergence. Its bulk spin RMS changes by only $0.0528\%$
over the last twenty records, but more than $99.999\%$ of the squared
final spin residual is orthogonal to an overall rescaling of the applied
spin field. Approximately $99\%$ of that residual power lies in the bulk.
The dominant remaining error concerns spatial shape.

For example, interpolated crossings of the staggered spin envelope near the
third and fourth walls shift from $39.3999$ and $55.1830$ at evaluation 60
to $39.5165$ and $55.2673$ at evaluation 80. Unequal shifts reveal texture
rearrangement. They are sub-rung profile diagnostics, not measured motion of
physical objects. The stopping diagnostic estimates a contraction near
$0.99059$, suggesting very slow relaxation if that behavior persists;
the extrapolation is neither a certified error bound nor a physical soft-mode
eigenvalue (\ev{zero}).

\subsection{Energy comparisons and what the bond-dimension check adds}

The relevant energy is the canonical variational functional, including the
mean-field double-counting contribution. If $L_{\mathrm{lin}}$ denotes
the expectation of the applied linear mean-field terms,
\begin{equation}
 E_{\mathrm{var}}=E_{\mathrm{eff}}+\mu N+
          (E_{\perp}-L_{\mathrm{lin}})
       =\avg{H_{\mathrm{lad}}}+E_{\perp}.
 \label{eq:energy}
\end{equation}
The density-targeted diagnostic is
$E_\star=E_{\mathrm{var}}+\mu(N_\star-N)$, a leading correction to the
target density. It does not remove curvature and optimization errors.
Effective-Hamiltonian eigenvalues alone cannot rank the branches
(\ev{methods}).

At $V=0$, the separate matched $\chi=400$ comparison gives a stripe endpoint
lower than the paired endpoint by $2.4336\times10^{-4}t$ per physical
site, or $0.03115t$ on the 128-site ladder. The September~6 working decision
treated the available evidence as sufficient to regard the advantage as
reasonably robust over the tested $\chi=200$ to 400 increase. Neither
original endpoint satisfies all formal acceptance gates, so the value
remains a terminal diagnostic difference rather than a fully certified
fixed-point ranking (\ev{chi}).

The newer raw $\chi=200$ trajectories instead end with a separation
$8.99\times10^{-5}t$ per site while the pairing lineage is still evolving.
This is not an independent estimate of the same equilibrium energy gap.
Similarly, an early low energy on an unconverged trajectory cannot be
selected as a lower-energy phase. During startup, the functional combines
the applied partner fields with current correlations; raw self-consistency
iteration is not guaranteed to minimize a single self-consistent state's
energy monotonically.

\subsection{Complete square and unfrustrated-cubic grids}
\label{sec:complete-grids}

The September~18 review extends the anchor comparison to
$t_0/t=1.0,1.2,1.4$ and $V/t=0,-0.2,-0.4$, with $U/t=8$,
$n=0.9375$, $t_\perp/t=0.1$, $L=64$ rungs and $\chi=200$.
The two starts are reciprocal 95\%/5\% mixtures of the reference stripe and
paired correlations, reconstructed with the target couplings. They evolve
with raw MF updates, without Anderson acceleration. All phase descriptions
below refer to finite-run physical correlations, not certified minima
(\ev{campaign}).

Figure~\ref{fig:square-cubic-phase} places the two phase diagrams side by
side with the same coordinate axes. After including both $(1.4,0)$
continuations, the coarse square grid has
seven stripe coordinates and two paired coordinates, $(1.4,-0.4)$ and
$(1.4,-0.2)$. Both seed families now reach the same qualitative basin at
each coordinate. At $(1.4,0)$ the stripe and pairing lineages contain
100 and 82 cumulative evaluations, respectively; the late pairing collapse
seen in the earlier history is therefore followed by a stripe endpoint.
The whole square grid contains 902 evaluations across 18 lineages, including
20 source jobs, with no formally accepted endpoint (\ev{square-complete}).

In contrast, all 18 unfrustrated-cubic starts reach essentially unpaired
CDW/SDW stripes after 60 evaluations. Physical leg-odd spin RMS ranges from
0.280 to 0.356, while the largest leg-pair RMS is $2.21\times10^{-11}$.
The dominant charge and spin modes are $m=4$ and 30, corresponding to
nominal 16-rung charge and 32-rung staggered-spin-envelope periods.
The two square paired coordinates are striped in cubic. Every cubic
endpoint passes its final energy-window, density, inner-DMRG and identity
checks; remaining field/profile and slow-mode checks prevent acceptance.
This is evidence of geometry-dependent nonlinear basin selection, not an
energetic ranking of different Hamiltonians or an exclusion of every other
cubic basin.

\begin{figure}[htbp]
 \centering
 \includegraphics[width=\textwidth]{\projectdocroot/reports/campaign_review_20260918/square_cubic_phase_diagrams.pdf}
 \caption{Preliminary square (left) and unfrustrated-cubic (right) phase
 diagrams. Both seeds reach the same observed family at every coordinate:
 seven stripe and two paired points on square, nine stripe points in cubic.
 S denotes CDW/SDW stripes and D opposite-sign leg/rung pairing.
 All endpoints remain unaccepted at $\chi=200$; no boundary is interpolated.
 Sources: \ev{square-complete} and \ev{campaign}.}
 \label{fig:square-cubic-phase}
\end{figure}

The full energy histories in Fig.~\ref{fig:square-cubic-energy} retain the
startup and delayed conversions. In particular, square $(1.2,-0.4)$ has an
extended pairing-seed energy plateau before its late decrease; square
$(1.4,0)$ evolves slowly into its continuation. Cubic energy histories
generally settle earlier. All plotted energies use the stored canonical
functional with target-density correction, and each panel keeps its own
vertical scale. The comparison is between numerical trajectories under
different geometries, not a ranking of their absolute energies as phases
of one Hamiltonian. Early off-self-consistent energy dips and late flat
energies do not establish stationary minima.

Figures~\ref{fig:square-cubic-spin} and \ref{fig:square-cubic-pairing}
show the complementary physical spin and pairing RMS histories. They use
the same definitions and logarithmic scales in both geometries, unlike the
older standalone square plots of MF-field amplitudes. Spin is the RMS of
the physical leg-odd spin on rungs 6--59; pairing is the physical leg-even
nearest-neighbor leg-pair RMS on bonds 6--58. Square spin decays while pairing
survives at $(1.4,-0.4)$ and $(1.4,-0.2)$; cubic instead retains strong spin
and loses pairing at both coordinates. At square $(1.2,-0.4)$ and $(1.4,0)$,
delayed pairing collapse accompanies increasing spin. The long transient in
these panels explains why a short, flat energy history can be misleading.
Tiny pairing fluctuations after collapse are numerical remnants, not
resolved renewed order.

\begin{landscape}
\begin{figure}[p]
 \centering
 \includegraphics[width=\linewidth]{\projectdocroot/reports/campaign_review_20260918/square_cubic_energy_grids.pdf}
 \caption{Full target-corrected canonical energy grids, square left and
 cubic unfrustrated right. Within each geometry, columns are $t_0=1,1.2,1.4$
 and rows are $V=0,-0.2,-0.4$. Solid blue and dashed orange denote stripe
 and pairing seeds. All 902 square and 1080 cubic saved evaluations are
 included. Dotted lines mark the square continuations after evaluations
 80 and 62; curves end at their own final evaluations. Matching coordinates
 share an iteration range, while energy y scales vary by panel.}
 \label{fig:square-cubic-energy}
\end{figure}
\begin{figure}[p]
 \centering
 \includegraphics[width=\linewidth]{\projectdocroot/reports/campaign_review_20260918/square_cubic_spin_grids.pdf}
 \caption{Physical leg-odd spin RMS histories, square left and cubic right,
 over rungs 6--59. All panels use a common logarithmic y scale. The point
 order, seed colors and continuation markers match
 Fig.~\ref{fig:square-cubic-energy}. These are physical correlations, not
 geometry-dependent Hartree-field amplitudes. Spin decays at the two square
 paired coordinates but remains strong throughout cubic.}
 \label{fig:square-cubic-spin}
\end{figure}
\begin{figure}[p]
 \centering
 \includegraphics[width=\linewidth]{\projectdocroot/reports/campaign_review_20260918/square_cubic_pairing_grids.pdf}
 \caption{Physical leg-pair RMS histories on bonds 6--58, square left and
 cubic right, with common logarithmic y scale. The layout matches
 Fig.~\ref{fig:square-cubic-spin}. Pairing survives at the two square paired
 points and decays to negligible amplitudes at every cubic point. The full
 histories retain delayed square pairing collapse, including the $V=0$
 continuation. None of the endpoints is formally accepted.}
 \label{fig:square-cubic-pairing}
\end{figure}
\end{landscape}

\subsection{Finer square cuts and the shape of the variational energy}
\label{sec:fine-energy}

The finer cuts add $V/t=-0.15,-0.10,-0.05$ at $t_0/t=1.4$ and
$t_0/t=1.25,1.30,1.35$ at $V/t=-0.4$, with 60 raw evaluations per seed.
Both seeds reach paired trajectories at four new points. Distinct stripe
and paired trajectories survive at $(1.25,-0.4)$ and $(1.4,-0.05)$.
Their signed endpoint energy differences, pairing seed minus stripe seed,
are $-3.189\times10^{-5}$ and $-5.952\times10^{-5}\,t$/site.
These exceed the sums of the corresponding final-ten energy ranges,
$1.574\times10^{-7}$ and $7.461\times10^{-8}\,t$/site, but recent drift
does not bound the remaining convergence error. All twelve endpoints remain
unaccepted (\ev{fine-cuts}).

Figure~\ref{fig:energy-cuts} shows the full stored variational energy along
both cuts, including the coarse endpoints and the latest $V=0$ continuations.
The quantity is the canonical functional with its field-dependent
double-counting terms and target-density correction,
\begin{equation}
 e_{\mathrm{target}}=E_{\mathrm{var}}/N_s+
 \mu(n_{\mathrm{target}}-n),\qquad N_s=2L=128.
\end{equation}
It is not the effective-Hamiltonian eigenvalue. No additional
field-independent perturbative offset has been introduced. On this scale
the two seed curves nearly overlap: the $V$ cut is almost linear, while
the $t_0$ cut has gradual curvature. Curves identify the initial seed,
not a persistent phase, because independent runs can change their outcome
as the parameter changes.

\begin{figure}[htbp]
 \centering
 \includegraphics[width=\textwidth]{\projectdocroot/reports/square_fine_cuts_20260918/variational_energy_cuts.pdf}
 \caption{Full canonical variational endpoint energies, corrected to the
 target density. Filled points are the new 60-step starts; open points use
 the coarse endpoints and latest continuations. Lines guide the eye between
 independent starts; they are not parameter-continuation branches.
 Total energy equals the displayed energy per site times 128.}
 \label{fig:energy-cuts}
\end{figure}

To resolve small shape changes, Fig.~\ref{fig:energy-shape} subtracts one
common linear chord from both seed curves on each cut. Its endpoints are
the mean seed energies at the two outer coordinates. This preserves both
seed gaps and changes in slope. The lower panels display adjacent secants
$\Delta e/\Delta p$, with $p=V$ or $t_0$, rather than a fitted derivative.
Along the pairing-seed $V$ curve the slopes are approximately 1.266231,
1.266144, 1.266212 and 1.261367. Thus the last interval, $-0.05\le V/t\le0$,
has a slope smaller by about 0.38\% than its predecessor. The stripe-seed
curve also bends downward in that interval. This small feature lies near
the change from paired to stripe outcomes and is compatible with a weak
first-order crossing, but does not resolve one.

For a first-order transition the minimizing energy can remain continuous
while its first derivative jumps. Five sampled coordinates do not
distinguish that kink from a continuous change of curvature. On the $t_0$
cut the pairing-seed secants, $-0.340515,-0.368857,-0.393241,-0.417749$,
change nearly regularly, with no obvious sharp cusp at this resolution.
The broad arch after linear subtraction is not itself a singularity.
Absence of a visible kink does not exclude a weak first-order transition.

\begin{figure}[p]
 \centering
 \includegraphics[width=\textwidth]{\projectdocroot/reports/square_fine_cuts_20260918/variational_energy_shape.pdf}
 \caption{Energy shape and adjacent-interval slopes. The same chord is
 subtracted from both seed curves within each cut; the two cuts have
 independent vertical scales. Bars on the secants are the sum of endpoint
 final-ten energy ranges divided by interval width, not statistical or
 convergence error bounds. In particular, the coarse $t_0=1.2$ pairing
 lineage still has an energy range $1.56\times10^{-4}\,t$/site.
 Polygonal corners between samples are not measured derivative jumps.}
 \label{fig:energy-shape}
\end{figure}

The two distinct textures at one point on each cut make metastability and
first order worth testing. However, these are not two stationary phase
branches traced through a crossing, nor parameter-sweep hysteresis.
The signed $E_p$ values at the new points were linearly interpolated between
the coarse endpoints as requested; $g=t_\perp^2/|E_p|$ therefore also changes
along each cut. The plotted slopes include this coupling dependence and
cannot be identified with the bare $V$ or rung-hopping conjugate observable
alone. Coarse and fine runs also have different stopping controls.

The weak spin at the paired $(1.4,-0.05)$ endpoint requires a spatial
qualification. Its usual spin RMS on rungs 6--59 rises 1.50\% over the last
ten evaluations to 0.005606, but 96.7\% of its spin-squared weight lies in
the outer 14 rungs at each end. Central 32-rung spin still decreases, while
leg-pair RMS remains 0.032991. This is not established bulk coexistence or
a demonstrated bulk return to stripes. At paired $(1.25,-0.4)$, the spin
remainder decreases by 52.1\% over the final ten evaluations.

Figure~\ref{fig:order-cuts} places spin and pairing alongside the energy
cuts. Coarse and fine points use the same physical definitions:
$M_i=(\langle S^z_{i,0}\rangle-\langle S^z_{i,1}\rangle)/2$, with spin RMS
over rungs 6--59, and symmetric nearest-leg anomalous pairing averaged
over the two legs before its RMS on bonds with left rungs 6--58. These
are not source-field norms or four-fermion pair--pair correlations.
Central-window spin exposes the end-weighted remainder at $(1.4,-0.05)$.
Pairing stays finite almost up to the sampled change of outcome, while
the stripe-seeded amplitudes change at a different sample. This motivates
testing competing finite-amplitude branches, but joining independent starts
does not measure a discontinuity, coexistence interval or hysteresis.
A logarithmic companion in the report resolves the small remainders.

\begin{figure}[p]
 \centering
 \includegraphics[width=\textwidth]{\projectdocroot/reports/square_fine_cuts_20260918/order_parameter_cuts.pdf}
 \caption{Spin and anomalous leg-pair RMS along both finer square cuts,
 using the same twenty endpoints as Fig.~\ref{fig:energy-cuts}. Filled
 markers denote fine starts; open markers denote coarse anchors with
 latest continuations. Dotted spin curves use central rungs 17--48.
 Lines join independent finite-iteration starts, not stationary phase
 branches. All endpoints remain unaccepted.}
 \label{fig:order-cuts}
\end{figure}

\clearpage
\subsection{Repulsive interaction on square: loss of pairing and magnetic defects}
\label{sec:positive-v}
\label{sec:positive-trellis}

All four square $(t_0/t,V/t)=(1.2,+0.2)$ starts now complete 60 raw
evaluations at $L=64$, $\chi=200$: the two 95\%/5\% reference mixtures
and intertwined patterns with initial charge wavelengths eight and sixteen.
The additional seeds place pairing near hole-rich regions and spin nodes,
testing whether the legacy frustrated-cubic texture survives on square.
Figure~\ref{fig:positive-histories} shows that all four lose anomalous pairing.
Stripe, pairing and period-16 starts end at spin RMS $0.23321$--$0.23358$
and leg-pair RMS below $8.3\times10^{-10}$. The period-eight start ends at
spin RMS $0.20700$ and leg-pair RMS $7.01\times10^{-9}$, with the latter
falling another 87.8\% over the last ten evaluations. None retains
appreciable anomalous coexistence (\ev{positive-v}).

\begin{figure}[p]
 \centering
 \includegraphics[width=\textwidth]{\projectdocroot/reports/square_positive_v_20260918/histories.pdf}
 \caption{All four square $(1.2,+0.2)$ histories. Full energy, spin RMS,
 logarithmic pairing RMS and late energy detail separate loss of pairing
 from persistence of different magnetic textures. Every start reaches the
 60-evaluation cap without acceptance.}
 \label{fig:positive-histories}
\end{figure}

The shorter-wavelength start has not merged with the other three.
Figure~\ref{fig:positive-profiles} shows four spin-envelope nodes in the
stripe, pairing and period-16 endpoints, near rungs 10, 25, 40 and 55.
Their charge/spin Fourier bins are $m=4/30$, giving nominal charge and
spin-envelope periods 16/32. The period-eight start evolves from eight
nodes at evaluation one to six near $10.16,23.21,26.81,38.30,41.89,54.91$.
Two close node pairs survive around the inner hole-rich regions. Charge
still has dominant $m=4$, while the spin spectrum has comparable
$m=31,30,28$ amplitudes $0.09264,0.08450,0.06761$. Its largest bin should
not be interpreted as a clean new spin wavelength. The spatial evidence
is consistent with incomplete coarsening or a trapped magnetic defect
pattern, rather than a surviving period-eight intertwined state.

\begin{figure}[p]
 \centering
 \includegraphics[width=0.76\textwidth,height=0.76\textheight,keepaspectratio]{\projectdocroot/reports/square_positive_v_20260918/terminal_profiles.pdf}
 \caption{Square positive-$V$ charge, staggered leg-odd spin and pairing
 profiles. Strong curves are evaluation 60, faint curves evaluation 51;
 dotted pairing curves are rung amplitudes. Only a global spin reversal
 aligns seeds. The period-eight start leaves two additional spin nodes.
 Pairing uses a $10^{-7}$ scale and is negligible in every endpoint;
 its residual spatial structure does not establish coexistence.}
 \label{fig:positive-profiles}
\end{figure}

The period-eight energy is $-0.332811625875\,t$/site, versus approximately
$-0.333744\,t$/site for the other three. Its excess over the stripe-seeded
endpoint is $9.3242\times10^{-4}\,t$/site; their summed final-ten energy
ranges are $8.96\times10^{-8}\,t$/site. This is a useful diagnostic
distinction, not a certified metastable branch energy. All four fail
global field/slow-mode and spatial charge, spin and exchange checks.
Maximum pointwise spin changes over the last ten evaluations remain
$0.00358$--$0.00655$, despite nearly constant RMS. Stripe and period-eight
starts also fail an inner-DMRG window; period sixteen fails the energy
window. Flat energy and a steady norm conceal moving walls.

The completed wavelength test strengthens the qualification on the earlier
physical interpretation: repulsive $V$ alone does not guarantee intertwining
across these geometries. Other square parameters or basins remain possible.
Loss of anomalous order also does not determine whether connected
pair--pair correlations survive in the stripes. The completed retrospective
measurements address that separate question in
Section~\ref{sec:pair-correlations}, without changing these acceptance flags.

\clearpage
\subsection{Trellis: spatial-cell dependence and alternating stripe relaxation}
\label{sec:trellis-results}

All four trellis comparisons now complete 60 raw cell sweeps at
$(t_0/t,V/t)=(1,0)$, $\tau_0/t=\tau_1/t=0.1$, density $15/16$ and
$\chi=200$, with $L=64$ per spatial ladder. Two seeds were tested in
the skew one-ladder cell and two in the rectangular $A/B$ cell.
All four remain unaccepted maximum-iteration endpoints. The contrast in
Fig.~\ref{fig:trellis-histories} is between spatial ansatzes: both one-ladder
starts are pairing dominated, while both two-ladder starts develop strong,
essentially unpaired stripes. The first one-ladder result must therefore
be qualified by this additional cell dependence (\ev{trellis}).

\begin{figure}[p]
 \centering
 \includegraphics[width=\textwidth]{\projectdocroot/reports/trellis_progress_20260918/histories.pdf}
 \caption{Complete trellis histories for both seeds in each spatial cell.
 Energy uses 128 sites for the skew one-ladder cell and 256 for the
 rectangular two-ladder cell. Spin and anomalous pairing use physical
 correlations and the same RMS windows as the square cuts. $A/B$ are
 spatial ladders, not successive iterations or temporal orbit phases.
 Vertical scales differ between cells.}
 \label{fig:trellis-histories}
\end{figure}

In the one-ladder cell, the stripe seed loses spin from $0.0627$ to
$1.215\times10^{-5}$ while leg-pair RMS rises from $0.00965$ to
$0.01796892$. The paired seed ends at spin RMS $1.368\times10^{-5}$
and leg-pair RMS $0.01796878$. Both have mean leg/rung amplitudes near
$+0.01796/-0.03292$: d-wave-like internal signs, not a spatial pair-density
wave. Their energies, $-0.518820336467$ and $-0.518820322856\,t$/site,
differ by $1.361\times10^{-8}\,t$/site. Maximum seed-to-seed full-profile
differences are $6.27\times10^{-6}$ in charge and $5.68\times10^{-7}$ in
leg pairing. This supports a common paired basin, although final-ten
energy ranges near $2.84\times10^{-7}\,t$/site and inner-DMRG/profile
failures prevent acceptance. Small spin fluctuations persist near $10^{-5}$.

Figure~\ref{fig:trellis-profiles} contrasts the one-ladder paired profiles
with the two-ladder stripes. Central charge standard deviation is about
$0.00157$ for one ladder, versus $0.0467$--$0.0470$ for two. End-sensitive
charge structure in the former does not establish bulk intertwining.
The latter retains strong central modulation, spin RMS
$0.22823$--$0.22876$ on both spatial ladders, and leg-pair RMS no larger
than $1.26\times10^{-8}$. Charge/spin bins $m=4/30$ correspond to nominal
periods 16/32. Charge minima and spin nodes shift between $A$ and $B$;
the latter is plotted at its physical offset of $-1/2$ rung.
Profiles are read directly from stored correlations: trellis Hartree
fields also contain normal-bond terms and cannot be inverted with the
square/cubic density formula.

\begin{landscape}
\begin{figure}[p]
 \centering
 \includegraphics[width=\linewidth,height=0.78\textheight,keepaspectratio]{\projectdocroot/reports/trellis_progress_20260918/profiles.pdf}
 \caption{Trellis profiles for all four starts and six spatial ladders.
 Solid curves are the final sweep, faint curves nine sweeps earlier, and
 dotted curves final rung pairing. Raw signs are retained. One-ladder
 spin and two-ladder pairing panels use tiny scales; neither establishes
 an additional ordered component. The two-ladder spatial displacement
 changes with sweep parity, as resolved in Fig.~\ref{fig:trellis-relaxation}.}
 \label{fig:trellis-profiles}
\end{figure}
\end{landscape}

Two-ladder nonacceptance is substantial, not marginal. Final global
relative residuals are $0.1845$--$0.2085$ against $10^{-4}$ allowed.
Both final inner-DMRG windows pass, but the outer MF map and its spatial
windows do not. Nearly constant spin RMS conceals the alternating wall
displacement in Fig.~\ref{fig:trellis-relaxation}. Successive full-cell
measured-field increments have cosine $-0.999962$ and $-0.999943$ for
the stripe and pairing starts, with norm ratios about $0.975$. Maximum
field changes are $0.02345/0.02620$ over one sweep but
$0.000754/0.000910$ over two. The iteration nearly reverses its previous
displacement while its amplitude slowly decreases. Two-sweep drift is
still resolved. This is consistent with a damped alternating numerical
mode, not an accepted period-two orbit or physical time-dependent order.
Modest linear damping is a possible future diagnostic of that negative
mode; no mixing change is made here.

\begin{figure}[p]
 \centering
 \includegraphics[width=\textwidth]{\projectdocroot/reports/trellis_progress_20260918/two_ladder_relaxation.pdf}
 \caption{Two-ladder relaxation. Upper panels compare measured fields
 separated by one and two sweeps, taking the maximum over all channels
 and both spatial ladders. Lower panels resolve staggered spin on $A$
 at sweeps 58, 59 and 60. Even sweeps nearly overlap while the odd sweep
 has displaced nodes. Large one-step changes and continuing two-step
 drift preclude stationarity.}
 \label{fig:trellis-relaxation}
\end{figure}

The two-ladder stripe- and pairing-seeded endpoint energies are
$-0.521125816792$ and $-0.521055560230\,t$/site. Their difference,
$7.026\times10^{-5}\,t$/site, is smaller than their summed final-ten
ranges of $3.445\times10^{-4}\,t$/site. To resolve the finite-cell energy
comparison, a September~20 audit places two identical copies of each saved
paired ladder state in the rectangular cell and recomputes all transverse
contractions there. Their target-corrected energies are
$-0.518817872689$ and $-0.518817859078\,t$/site, only
$2.46378\times10^{-6}\,t$/site above the original paired values.
Both striped trials thus remain $0.002238$--$0.002308\,t$/site lower
under the same finite-cell functional. Density corrections are below
$8.33\times10^{-6}\,t$/site and do not determine the ordering.
The open-end embedding difference is consequently much too small to
explain the observed energy advantage for these specified trials.
Self-consistency is not needed to compare their expectation values.
The large residuals still prevent identifying optimized stationary
branches or a unique global phase minimum. Both seeds agree on the
qualitative order within each cell, and a lower-energy striped trial
has now been established in a common cell.

A remaining physical restriction is the transverse stripe period. After
correcting the half-rung origin of $B$, the central charge harmonic at
$q_x=\pi/8$ has a relative $B$-minus-$A$ phase near $126$--$131$ degrees.
Repeating $A/B$ alternates this phase advance and does not establish
uniformly tilted, diagonal walls across many ladders. For an ideal charge
profile $\delta n_m(x)=\rho\cos(q_xx+m\theta)$, a rectangular $n$-ladder
cell requires $n\theta=2\pi p$. Equivalently, for a rigid shift of $d$
rungs per ladder, $nd$ must close the longitudinal charge period and the
full spin texture. Nominal charge/spin periods 16/32 with $d=2$ would,
for example, require sixteen ladders to repeat both without an additional
spin transformation. This is an illustrative closure condition, not a
measured tilt or proposed run size. The actual cell translation, trellis
shear and leg basis must enter any such construction. Larger transverse
cells or suitable translated repetitions could test oblique stripes;
the current two-cell comparison does not determine their optimum.

\clearpage
\subsection{Complete pair correlations: local survival and spatial suppression}
\label{sec:pair-correlations}

The September~22 synchronized measurements cover all 56 retrospective
branches (58 spatial MPSs), plus eight diagnostics from the four newer
square A/B runs. Receipts, full-source identifiers, compact-state hashes,
model/configuration provenance, connected subtraction, Hermiticity and
Gram positivity have been checked (\ev{pair-correlations}). All retrospective
branches remain unaccepted maximum-iteration snapshots. The accepted square
A/B states are discussed separately below.

For this section, use the stored unnormalized singlet
$D_{ab}=\sqrt2\Delta_{ab}$, with $\Delta$ defined in
Eq.~\eqref{eq:pair-op}, and define
\begin{equation}
 \widetilde P_{ab,cd}=\avg{D^\dagger_{ab}D_{cd}},\qquad
 \widetilde P^{\rm c}_{ab,cd}=\widetilde P_{ab,cd}
       -\avg{D_{ab}}^*\avg{D_{cd}}.
 \label{eq:connected-pair-measured}
\end{equation}
All new pair-correlation numbers and figures use this stored convention;
divide by two for Eq.~\eqref{eq:pair-op}'s normalized bond correlations.
The complete basis has 318 onsite, rung and nearest-neighbor leg operators,
including cross-channel entries. The older isolated reference stores only
within-class Hermitian-field covariances. For disjoint bonds in that real,
number-conserving state, its non-onsite field-metric factor of $1/2$
converts the sum of addition/removal correlations to
$\operatorname{Re}\widetilde P=\operatorname{Re}\widetilde P^{\rm c}$.
Contact and overlapping-bond entries are excluded; no missing isolated
rung--leg cross-correlations are inferred.

Figure~\ref{fig:pair-decay} averages correlation magnitudes over both
directions with rung bonds in 9--56 and leg-bond starts in 9--55, keeping
both endpoints inside that bulk window. The isolated reference uses
$\chi=1200$, whereas coupled states use $\chi=200$; this comparison is not a
matched truncation-error study. The square stripe at $(t_0,V)=(1.4,0)$
retains a mean connected rung magnitude $1.51\times10^{-2}$ at separations
2--4, compared with $2.03\times10^{-2}$ in the isolated reference. At
16--24 rungs the corresponding values are $9.81\times10^{-5}$ and
$5.22\times10^{-3}$. Local pairing survives, but its spatial persistence
is strongly reduced in these saved states. Both stripe seeds support the
same qualitative comparison.

\begin{figure}[htbp]\centering
\includegraphics[width=\textwidth]{\projectdocroot/reports/pair_correlations_20260922/pair_decay_comparison.pdf}
\caption{Full and connected rung-pair magnitudes and normalized connected
decay. Figure labels $P$ denote $\widetilde P$ in
Eq.~\eqref{eq:connected-pair-measured}. All coupled curves here are
unaccepted snapshots; the isolated curve has a larger bond dimension.
Square/cubic use $t_0=1.4$; trellis uses $t_0=1,V=0$.
Very small far-distance tails have no independent truncation-error estimate.}
\label{fig:pair-decay}
\end{figure}

At matched intraladder parameters $(1.4,-0.4)$ and $\chi=200$, square
paired and cubic striped endpoints have comparable short-distance
connected magnitudes ($1.18\times10^{-2}$ versus approximately
$1.01$--$1.03\times10^{-2}$). At 16--24 rungs, the square value is
$3.04\times10^{-4}$ and cubic gives $1.29$--$1.39\times10^{-7}$, a
factor of roughly $2.2$--$2.3\times10^3$. The trellis comparison at
$(1,0)$ is similarly selective: short-distance connected amplitudes in
the paired one-ladder and striped two-ladder endpoints differ by only
a few percent, while the long-distance value falls from
$4.64\times10^{-5}$ to approximately $6.6\times10^{-7}$.
Figure~\ref{fig:pair-matched} includes both seeds and spatial ladders.
These are endpoint and embedding comparisons, not a causal intervention
or a stationary phase ranking.

\begin{figure}[htbp]\centering
\includegraphics[width=\textwidth]{\projectdocroot/reports/pair_correlations_20260922/matched_pair_comparisons.pdf}
\caption{Matched intraladder-parameter comparisons of full and connected
pair correlations at $\chi=200$. Both initial families and each spatial
ladder are shown. The states retain their original unaccepted status.}
\label{fig:pair-matched}
\end{figure}

The square paired state's full long-distance rung correlation is
$1.17\times10^{-2}$: approximately 97.4\% of its signed mean is the
disconnected anomalous product. Full and connected correlations instead
coincide to the displayed precision in the representative stripe states.
A full-correlation plateau therefore cannot be identified with slowly
decaying connected fluctuations.

Relative signs also survive. In the representative square paired, square
striped and cubic striped states, rung--rung and leg--leg averages are
positive, while rung--leg averages are negative. Across all retrospective
MPSs, the fraction of negative connected rung--leg0 entries at bulk
separations 2--8 is 96.3--100\% on square, 87.2--100\% on cubic and
100\% on trellis. This supports a surviving local rung/leg sign structure,
commonly termed $d$-wave-like, without establishing a condensate or PDW.
Trellis leg--leg correlations also change sign with separation, so a
single uniform form factor would overstate the evidence. Signed channel
curves and maps are retained in the report (\ev{pair-correlations}).

The surviving short-distance pairing is spatially organized. Define
$A(i)=\sum_{r=2}^4(|\widetilde P^{\rm c}(i,i+r)|+
|\widetilde P^{\rm c}(i,i-r)|)/6$ for rung-pair centers 13--52, whose six
neighbors all lie inside the common bulk window. The mean $A$ in the
highest hole-density quartile exceeds the lowest quartile by factors
5.85--6.68 for square $V=0$, 19.39--19.59 for cubic $(1.4,-0.4)$, and
5.29--5.36 across both seeds/ladders of striped trellis. Interior hole
maxima lie within half a rung of the discrete staggered-spin sign crossings
(Fig.~\ref{fig:pair-walls}). These within-texture associations are
descriptive; they do not establish that stripes generate an attractive
interaction or a dynamical pairing mechanism.

\begin{figure}[htbp]\centering
\includegraphics[width=\textwidth]{\projectdocroot/reports/pair_correlations_20260922/stripe_pairing_profiles.pdf}
\caption{Short-distance pair correlations concentrate near hole-rich
antiphase magnetic walls in representative snapshots. The bottom panels
show $A(i)$; the middle panels show staggered leg-odd spin. Outer shading
excludes eight rungs and the pair-strength centers use rungs 13--52.
Dotted lines mark interior hole maxima; no independent-sample statistics
or causal mechanism is inferred.}
\label{fig:pair-walls}
\end{figure}

Reference sensitivity prevents a universal decay-exponent claim. Fits of
$\log|\widetilde P^{\rm c}(i,i\pm r)|$ at reference rungs 16, 24 and 32,
with windows 4--12, 8--20 and 12--28 clipped to the bulk, give effective
power slopes 0.39--1.36 for the isolated reference, 1.55--4.12 for square
paired, $-0.50$--6.74 for square striped, 0.37--15.55 for cubic striped
and 1.92--13.64 for the displayed trellis stripe. Fits require six points
above $10^{-12}$; this analysis threshold is not an error bound. The
negative slope is a local modulation effect, not asymptotic growth.
Alternate edge cuts and same-window exponential fits are also supplied.
This sensitivity reinforces the need for reference, boundary, length and
truncation controls emphasized by \citet{Shen2023}. Neither the old isolated
fit near 0.76 nor a selected new window should be treated as a universal
exponent without those controls.

\subsection{Square A/B: all four attractive runs reach paired fixed points}
\label{sec:square-ab-complete}

The square spatial extension at $t_0=1.4$ now has complete states and
measurements for both seeds at $V=-0.4$ and $-0.2$
(\ev{square-ab-complete}). The cell contains two 64-rung MPSs with
unshifted reciprocal square bonds and separate mean densities $15/16$.
The stripe component of ladder B's initial condition is displaced by
eight rungs. Unlike the rectangular trellis result, all four square
trajectories pass the fixed-point gates and lose their seeded magnetic
texture (Fig.~\ref{fig:square-ab-complete}).

\begin{table}[htbp]\centering
\caption{Accepted square A/B endpoints, $t_0=1.4$, $\chi=200$.
The pair RMS uses the historical physical symmetrized nearest-neighbor
leg anomalous expectation on bonds 6--58, not $\widetilde P$.}
\label{tab:square-ab-complete}
\begin{tabular}{r l r r r}\toprule
$V$ & Seed & Sweeps & Corrected $E/(Nt)$ & Leg-pair RMS\\\midrule
$-0.4$ & Stripe & 40 & $-1.167404800596$ & $0.0352251$\\
$-0.4$ & Pairing & 40 & $-1.167404800545$ & $0.0352251$\\
$-0.2$ & Stripe & 55 & $-0.914117842951$ & $0.0335678$\\
$-0.2$ & Pairing & 44 & $-0.914117843004$ & $0.0335678$\\\bottomrule
\end{tabular}
\end{table}

The two seed energies at each $V$ agree within $5.3\times10^{-11}\,t$/site.
Bulk spin RMS is below $5.4\times10^{-7}$ at $V=-0.4$ and
$1.15\times10^{-6}$ at $V=-0.2$. Maximum A/B leg-pair profile differences
are below $1.1\times10^{-6}$; charge differences reach $1.13\times10^{-5}$.
The independent audit checks raw-update identities, source and seed hashes,
channel/density/energy windows, inner-DMRG and slow-mode criteria, and
reconstructs canonical energies with 256 sites per cell.

\begin{figure}[htbp]\centering
\includegraphics[width=\textwidth]{\projectdocroot/reports/pair_correlations_20260922/square_AB_complete.pdf}
\caption{All four square A/B trajectories reach accepted paired fixed
points. Both seed families and both spatial ladders are shown. The seed
magnetization decays; the endpoint leg-pair profiles nearly coincide.}
\label{fig:square-ab-complete}
\end{figure}

The new pair correlations also agree between seeds and ladders and reproduce
the earlier paired one-ladder snapshots (Fig.~\ref{fig:square-ab-pairs}).
Connected rung magnitudes over 16--24 rungs are about $3.036\times10^{-4}$
and $2.475\times10^{-4}$ at $V=-0.4$ and $-0.2$; the corresponding full
values are $1.166\times10^{-2}$ and $1.030\times10^{-2}$. These results
support survival of the paired state in this tested two-ladder square cell.
They do not prove linear stability against every transverse mode, global
minimality or thermodynamic superconductivity. The diagnostics remain
intraladder: A/B mean-field coupling does not supply interladder quantum
entanglement or a phase-stiffness measurement.

\begin{figure}[htbp]\centering
\includegraphics[width=\textwidth]{\projectdocroot/reports/pair_correlations_20260922/square_AB_pair_correlations.pdf}
\caption{Connected rung-pair correlations in accepted square A/B states
and earlier one-ladder snapshots. Agreement holds for both seeds and both
attractions. The old one-ladder acceptance flags remain unchanged.}
\label{fig:square-ab-pairs}
\end{figure}

\clearpage
\subsection{Scope of the emerging physical interpretation}

The isolated-ladder and coupled-array observations separate pair formation
from ordered-state selection. A spin-gapped ladder with slowly decaying pair
correlations can still feed a static stripe state once transverse normal
fields are allowed to develop. At the attractive anchor, both tested starts
instead approach the same pairing-dominated solution. The current evidence
therefore supports a model-specific competition between collective orders
within a system that already possesses local pairing tendencies.

This result shares an ingredient with stripe-centered pairing mechanisms:
correlated low-dimensional subsystems can support pairing before collective
coherence is established. At the level of the current static textures,
however, the anchor comparison resembles competition between stripe order
and superconducting order. It neither establishes that fluctuating stripes
are necessary for pairing nor excludes their possible beneficial role.
The approximation treats quantum fluctuations within the ladder, but does
not resolve the full collective stripe and phase fluctuations of a plane.

The completed four-fermion measurements now resolve part of the earlier
open question: the measured stripe endpoints retain local singlet-pair
correlations near their hole-rich walls, while longer-distance correlations
are strongly suppressed. This supports a distinction between local pairing
and spatially extended order, with direct competition evident at longer
distances. It does not isolate phase fluctuations as the cause, establish
fluctuating superconductivity or certify the unaccepted stripe endpoints as
stationary phases. The newly accepted square A/B paired states strengthen
the paired side of the comparison within that specific spatial extension.

Attraction changes more than one ingredient of the coupled problem.
The prefactor $g=\tp^2/|\Ep|$ decreases when $|\Ep|$ increases at fixed
$\tp$, even while the ladder's pair response can become stronger.
This motivates separating the change in the bare ladder from the change
in the coupling kernel across a $V$ scan. Recent MPS+MF work on a different,
mixed-dimensional $t$--$J$ model likewise emphasizes the properties of the
one-dimensional building blocks and the requirements for an ordered
array~\citep{Kohler2026}. That work is retained as an August~2026 preprint,
not a result for the present Hamiltonian.

\subsection{Controls needed before stronger claims}

The pair-breaking expansion needs a careful energy-scale comparison.
The backbone gives $\tp/|\Ep|\simeq0.683$ and
$\tp/\Delta_s\simeq0.652$ at $\tp/t=0.1$; these ratios are not
parametrically small. The quantity
\begin{equation}
 \Delta_{c,2}=\frac{E(N+2)+E(N-2)-2E(N)}{2}
\end{equation}
instead measures even-particle charge curvature. A gapless charge mode can
coexist with a finite spin or pair-breaking scale, so
$\tp/\Delta_{c,2}>1$ does not by itself invalidate the projection.
This corrects the stronger interpretation in the earliest bare-ladder
report (\ev{review}). The registry and backbone also use slightly different
binding energies, differing by about $0.0188\%$ in the inferred coupling;
quantitative response/SCF comparisons should declare a common kernel.

The remaining high-value controls are:
\begin{enumerate}
 \item Complete self-consistency at the representative stripe point and
 compare matched canonical energies only where the states and provenance
 support it. A converged fixed point and a global minimum remain different
 claims.
 \item Test selected lengths, bond dimensions, and stripe wavelengths.
 The $L=96$ and 128 seeds are prepared, but no result from them is used here.
 Their planned fixed-$L=64$ binding denominator isolates one finite-size
 question; it is not a new measurement of the larger ladder's binding
 (\ev{length}).
 \item Check sensitivity to the mean-field range cutoff, currently
 $r_{\mathrm{range}}=4$ in the representative comparisons, and to weaker
 transverse hopping. These are model/approximation controls, separate from
 inner-DMRG convergence.
 \item Complete amplitude and basis checks for the response pilot and measure
 response about a self-consistent normal reference at fixed density.
 \item Separate connected fluctuations from disconnected order contributions.
 Subtracting $\avg{O_i}\avg{O_j}$ removes the elastic mean-field order
 contribution; both parts are needed to interpret structure factors.
 Extend the completed bond-class and reference-position comparisons to
 matched larger lengths and bond dimensions; the current effective slopes
 have substantial window dependence.
\end{enumerate}

The completed square and unfrustrated-cubic grids now supply contrasting
finite-run phase patterns. The finer cuts identify competing textures but
do not resolve the order of the transition. The final September~19 review
contains 38 completed 60-step run histories, or 2280 cell updates and
2400 individual ladder solves, all unaccepted;
the earlier coarse square grid is a separate data set. No universal effect
of $V$, certified ground-state phase diagram, transition
temperature, phase stiffness, pair-density-wave state, or dynamical magnetic
spectrum is claimed. The present contribution is a developing account of
how local pairing tendencies and competing static orders relate in a
specified coupled-ladder construction.

\appendix
\section{Local evidence and drafting notes}
\label{sec:evidence}

\subsection{Sources for project-specific statements}

The following sources are repository documents and analysis artifacts,
not external publications. Paths are relative to \path{ladder_mps_mft/docs}.
The linked reports retain the underlying histories, definitions, and
provenance checks. Their dated conclusions must be read together: later
history analyses supersede the early suggestion of a unique paired
$V=0$ basin.

\begin{enumerate}
 \item\label{ev:methods}
 \textbf{Model, fields, and energy.}
 \href{run:../METHODS_NOTES.tex}{Living methods notes},
 \href{run:../VARIATIONAL_FUNCTIONAL.md}{variational functional}, and
 \href{run:../ARCHITECTURE.md}{architecture}.
 \path{METHODS_NOTES.tex};
 \path{VARIATIONAL_FUNCTIONAL.md}; \path{ARCHITECTURE.md}.

 \item\label{ev:geometry}
 \textbf{Literature and equation-level geometry correspondence.}
 \href{run:../literature/literature_review.pdf}{Literature review} and
 \href{run:../literature/SOURCE_NOTES.md}{source notes}, including the
 September~8 comparison with the 2023 benchmark.
 \path{literature/literature_review.tex};
 \path{literature/SOURCE_NOTES.md}.

 \item\label{ev:bare}
 \textbf{Isolated-ladder backbone and covariance survey, September~2.}
 \href{run:../reports/bare_stage1_t014_v0_20260902/REPORT.md}{Stage 1 report}.
 \path{reports/bare_stage1_t014_v0_20260902/REPORT.md}.

 \item\label{ev:response}
 \textbf{Projected response discovery, September~3.}
 \href{run:../reports/bare_stage2_t014_v0_20260903/report.html}{Stage 2 report},
 with eigenvalues and reference-field summaries.
 \path{reports/bare_stage2_t014_v0_20260903/}.

 \item\label{ev:review}
 \textbf{Cross-project scientific review, September~4.}
 \href{run:../reports/systematic_review_20260904/REVIEW.md}{Review}.
 Source for the response limitations, charge-gap correction, and proposed
 approximation controls; its campaign statuses are historical.
 \path{reports/systematic_review_20260904/REVIEW.md}.

 \item\label{ev:chi}
 \textbf{Bond-dimension comparison, September~5--6.}
 \href{run:../reports/chi400_comparison_20260905/ANALYSIS.md}{Analysis},
 including the explicitly qualified September~6 energetic interpretation.
 \path{reports/chi400_comparison_20260905/ANALYSIS.md}.

 \item\label{ev:basins}
 \textbf{History-based revision of basin interpretation, September~8.}
 \href{run:../reports/square_grid_20260908/BASIN_ASSESSMENT.md}{Basin assessment}.
 \path{reports/square_grid_20260908/BASIN_ASSESSMENT.md}.

 \item\label{ev:grid}
 \textbf{Square coverage and the flagged $(1.0,-0.4)$ endpoint, September~8.}
 \href{run:../reports/square_grid_20260908/README.md}{Grid description} and
 \href{run:../reports/square_grid_20260908/ANALYSIS.md}{divergence analysis}.
 \path{reports/square_grid_20260908/}.

 \item\label{ev:attractive}
 \textbf{Attractive two-family anchors, September~10.}
 \href{run:../reports/two_basin_vm04_20260910/README.md}{Analysis};
 complementary retrospective gate qualification in the
 \href{run:../reports/two_basin_remainder_20260910/README.md}{remainder contract}.
 \path{reports/two_basin_vm04_20260910/README.md}.

 \item\label{ev:zero}
 \textbf{Zero-attraction anchors and slow texture relaxation, September~12--13.}
 \href{run:../reports/two_basin_v000_20260912/README.md}{Analysis and follow-up}.
 Figure~\ref{fig:zero} is reused without modifying the underlying data.
 \path{reports/two_basin_v000_20260912/README.md}.

 \item\label{ev:campaign}
 \textbf{Combined synchronized campaign review, updated September~19.}
 \href{run:../reports/campaign_review_20260918/README.md}{Combined report} and
 \href{run:../reports/cubic_two_basin_grid_20260918/README.md}{complete cubic grid}.
 Source inventories, physical profiles, gate checks and accounting coverage
 are retained with the detailed reports.
 The side-by-side figures use
 \href{run:../reports/campaign_review_20260918/square_cubic_histories.csv}{matched physical histories}
 and their \href{run:../reports/campaign_review_20260918/square_cubic_comparison_sources.json}{source inventory}.

 \item\label{ev:square-complete}
 \textbf{Complete coarse square grid and continuations.}
 \href{run:../reports/two_basin_grid_20260915/README.md}{Updated square report},
 including both $(1.4,0)$ continuations and the historical parent records.

 \item\label{ev:fine-cuts}
 \textbf{Finer square cuts and full energy curves, September~18.}
 \href{run:../reports/square_fine_cuts_20260918/README.md}{Cut report},
 \href{run:../reports/square_fine_cuts_20260918/variational_energy_cuts.csv}{endpoint energy table},
 and \href{run:../reports/square_fine_cuts_20260918/variational_energy_analysis.json}{energy audit}.
 Twenty endpoint sources include the twelve fine runs and eight coarse anchors.
 The \href{run:../reports/square_fine_cuts_20260918/order_parameter_analysis.json}{order-parameter audit}
 uses the same sources and consistent physical RMS windows.

 \item\label{ev:positive-v}
 \textbf{Square repulsive-interaction test, completed September~19.}
 \href{run:../reports/square_positive_v_20260918/README.md}{Four-seed comparison};
 all four spatial endpoints, with spin-node and energy diagnostics.

 \item\label{ev:trellis}
 \textbf{All four trellis endpoints and alternating relaxation, September~19.}
 \href{run:../reports/trellis_progress_20260918/README.md}{Trellis report} and
 \href{run:../TRELLIS_MEAN_FIELD.md}{spatial-cell method contract}.

 \item\label{ev:pair-correlations}
 \textbf{Complete retrospective pair correlations, September~22.}
 \href{run:../reports/pair_correlations_20260922/correlation_report.pdf}{Correlation report},
 with full/connected and signed channel figures, source hashes, complete
 coverage, reference-window fits, bulk-cut sensitivity and stripe-wall metrics.
 \path{reports/pair_correlations_20260922/}.

 \item\label{ev:square-ab-complete}
 \textbf{All four accepted square A/B results, September~22.}
 The same report's completed square section and
 \href{run:../reports/pair_correlations_20260922/square_AB_results.json}{gate/energy audit}
 supersede the September~21 first-result-only status.
 The original \href{run:../reports/square_two_ladder_20260920/README.md}{preparation contract}
 retains the model, seeds and spatial-cell definition.

 \item\label{ev:length}
 \textbf{Prepared finite-size seeds, September~6.}
 \href{run:../reports/finite_size_seeds_20260906/README.md}{Seed description}.
 Preparation evidence only; no larger-length simulation result is inferred.
 \path{reports/finite_size_seeds_20260906/README.md}.
\end{enumerate}

\subsection{Reference provenance and manuscript development}

The draft reuses stable keys from \path{../literature/references.bib}.
The initial five additional sources in \path{additional_references.bib} are
Tranquada et al.\ (2004), Vojta and Ulbricht (2004), Li et al.\ (2007),
Berg et al.\ (2007), and Missiaen et al.\ (2025). Their primary arXiv and
publisher records were checked on 15 September 2026. These additions
support the discussion of magnetic spectra, planar coherence, and the
stripe--pseudogap connection. Eleven further sources, checked on
18 September 2026, support the ladder-material, doping, optical-parameter,
charge-order, NMR, neutron, and RIXS discussion in
Section~\ref{sec:ladder-materials}; reading depths and primary links are
recorded in \path{SOURCE_NOTES.md}. The existing review's reading-depth limits
continue to apply to inherited references; this draft is not an exhaustive
literature or novelty assessment.

The introduction can be shortened for a target journal after the principal
result and convergence controls are settled. The current results collection
should retain its date and evidence qualifications until replaced by
converged comparisons. A future manuscript should decide whether its central
claim is interaction-driven order selection, geometry dependence, or the
relation between a dressed normal-state response and the nonlinear ordered
solutions. Those are candidate emphases, not additional results.

\clearpage
\begingroup
\setlength{\bibsep}{0.20em}
\bibliographystyle{unsrtnat}
\bibliography{\projectdocroot/literature/references,\manuscriptroot/additional_references}
\endgroup
\end{document}
